\documentclass[a4paper, 10pt, landscape]{scrartcl}

% ============================================================
% PRÄAMBEL & LAYOUT (Design Herr Dayi - Comic-Vorlage)
% ============================================================
\usepackage[a4paper, landscape, top=0.9cm, bottom=0.9cm, left=1.2cm, right=1.2cm]{geometry}
\usepackage[ngerman]{babel}
\usepackage{libertinus}
\usepackage{xcolor}
\usepackage[most]{tcolorbox}
\usepackage{tabularx}
\usepackage{tikz}

\pagestyle{empty}

% Farbdefinitionen
\definecolor{primaryBlue}{RGB}{20, 50, 100}
\definecolor{lightBlue}{RGB}{245, 248, 253}
\definecolor{accentOrange}{RGB}{220, 100, 30}
\definecolor{panelBorder}{RGB}{180, 195, 215}

\begin{document}

% === KOPFZEILE ===
\noindent
\begin{minipage}[t]{0.32\textwidth}
    {\small\textbf{Philosophie Q1} \enskip$\cdot$\enskip Herr Dayi}
\end{minipage}%
\begin{minipage}[t]{0.36\textwidth}
    \centering{\small\textbf{Anthropologie \& Kultur (IF 3)}}
\end{minipage}%
\begin{minipage}[t]{0.32\textwidth}
    \flushright{\small Name: \rule{3.2cm}{0.4pt} \enskip Datum: \rule{1.8cm}{0.4pt}}
\end{minipage}

\vspace{0.12cm}
{\color{primaryBlue}\hrule height 1pt}
\vspace{0.18cm}

% === TITEL & ARBEITSAUFTRAG ===
\noindent
\begin{tcolorbox}[
    enhanced,
    colback=lightBlue,
    colframe=primaryBlue,
    boxrule=0.8pt,
    arc=1.5mm,
    top=4pt, bottom=4pt, left=8pt, right=8pt
]
\noindent
{\large\bfseries\color{primaryBlue} Comic-Strip: Die Entfremdung des Menschen -- Von Chaplin bis heute}\\[2pt]
\footnotesize\sffamily
\textbf{Arbeitsauftrag:} Stellt Karl Marx' Diagnose der entfremdeten Arbeit in einem \textbf{4-teiligen Comic-Strip} dar (einfache Strichmännchen genügen vollkommen!). Zeichnet die Verwandlung des Menschen vom freien Schöpfer zum fremdbestimmten Rädchen im Getriebe. Nutzt \textbf{Sprech- und Denkblasen}, Symbole und pointierte Alltagsszenen.
\end{tcolorbox}

\vspace{0.18cm}

% === 4 COMIC PANELS (NEBENEINANDER) ===
\noindent
\begin{minipage}[t]{0.238\textwidth}
    % PANEL 1
    \begin{tcolorbox}[
        enhanced,
        colback=white,
        colframe=primaryBlue,
        boxrule=1.2pt,
        arc=2mm,
        height=13.6cm,
        top=5pt, bottom=5pt, left=6pt, right=6pt,
        title={\footnotesize\bfseries\color{white} 1. Das Gattungswesen (Ideal)},
        boxed title style={colback=primaryBlue, arc=1mm, boxrule=0pt}
    ]
        \scriptsize
        \textit{\glqq[\dots] während der Mensch frei produziert [\dots] nach den Gesetzen der Schönheit.\grqq}\\[3pt]
        {\tiny\sffamily\color{primaryBlue}\textbf{Fokus:} Freie Gestaltung, Kunst, Sinn, Freude am Erschaffen}
        \vspace{4pt}
        {\color{primaryBlue!30}\hrule height 0.8pt}
    \end{tcolorbox}
\end{minipage}%
\hfill
\begin{minipage}[t]{0.238\textwidth}
    % PANEL 2
    \begin{tcolorbox}[
        enhanced,
        colback=white,
        colframe=primaryBlue,
        boxrule=1.2pt,
        arc=2mm,
        height=13.6cm,
        top=5pt, bottom=5pt, left=6pt, right=6pt,
        title={\footnotesize\bfseries\color{white} 2. Herrschaft des Produkts},
        boxed title style={colback=primaryBlue, arc=1mm, boxrule=0pt}
    ]
        \scriptsize
        \textit{\glqq Was das Produkt seiner Arbeit ist, ist er nicht. [\dots] Herrschaft des Kapitals.\grqq}\\[3pt]
        {\tiny\sffamily\color{primaryBlue}\textbf{Fokus:} Die Maschine / das Ding beherrscht den Arbeiter}
        \vspace{4pt}
        {\color{primaryBlue!30}\hrule height 0.8pt}
    \end{tcolorbox}
\end{minipage}%
\hfill
\begin{minipage}[t]{0.238\textwidth}
    % PANEL 3
    \begin{tcolorbox}[
        enhanced,
        colback=white,
        colframe=primaryBlue,
        boxrule=1.2pt,
        arc=2mm,
        height=13.6cm,
        top=5pt, bottom=5pt, left=6pt, right=6pt,
        title={\footnotesize\bfseries\color{white} 3. Zwangsarbeit \& Kasteiung},
        boxed title style={colback=primaryBlue, arc=1mm, boxrule=0pt}
    ]
        \scriptsize
        \textit{\glqq In der Arbeit außer sich [\dots] keine freie Energie, sondern Kasteiung der Physis.\grqq}\\[3pt]
        {\tiny\sffamily\color{primaryBlue}\textbf{Fokus:} Monotonie, Stechuhr, Zwang, Erschöpfung}
        \vspace{4pt}
        {\color{primaryBlue!30}\hrule height 0.8pt}
    \end{tcolorbox}
\end{minipage}%
\hfill
\begin{minipage}[t]{0.238\textwidth}
    % PANEL 4
    \begin{tcolorbox}[
        enhanced,
        colback=white,
        colframe=primaryBlue,
        boxrule=1.2pt,
        arc=2mm,
        height=13.6cm,
        top=5pt, bottom=5pt, left=6pt, right=6pt,
        title={\footnotesize\bfseries\color{white} 4. Das Paradoxon (Feierabend)},
        boxed title style={colback=primaryBlue, arc=1mm, boxrule=0pt}
    ]
        \scriptsize
        \textit{\glqq In seinen tierischen Funktionen freitätig und in menschlichen nur mehr als Tier.\grqq}\\[3pt]
        {\tiny\sffamily\color{primaryBlue}\textbf{Fokus:} Schein-Freiheit beim Essen, Trinken, Schlafen}
        \vspace{4pt}
        {\color{primaryBlue!30}\hrule height 0.8pt}
    \end{tcolorbox}
\end{minipage}

\vspace{0.22cm}

% === SCHLUSSZEILE: TITEL / POINTE DES COMICS ===
\noindent
\begin{tcolorbox}[
    enhanced,
    colback=white,
    colframe=primaryBlue,
    boxrule=0.8pt,
    arc=1mm,
    top=4pt, bottom=4pt, left=8pt, right=8pt
]
\small\sffamily
\textbf{\color{primaryBlue}Titel oder Schlusspointe deines Comic-Strips:} \rule{19.5cm}{0.4pt}
\end{tcolorbox}

\end{document}
